\documentclass[11pt,a4paper]{article}

\usepackage[utf8]{inputenc}
\usepackage[T1]{fontenc}
\usepackage{lmodern}
\usepackage[margin=0.82in]{geometry}
\usepackage[table]{xcolor}
\usepackage{array}
\usepackage{tabularx}
\usepackage{longtable}
\usepackage{booktabs}
\usepackage{enumitem}
\usepackage{titlesec}
\usepackage{fancyhdr}
\usepackage{lastpage}
\usepackage{hyperref}
\usepackage{multicol}

\definecolor{TextInk}{HTML}{233142}
\definecolor{HeadingBlue}{HTML}{1D3557}
\definecolor{AccentTeal}{HTML}{2A7F92}
\definecolor{AccentGold}{HTML}{C58B4E}
\definecolor{AccentGreen}{HTML}{4F7F52}
\definecolor{RuleGray}{HTML}{D6E0E5}
\definecolor{TableStripe}{HTML}{F3F8FA}
\definecolor{SoftTint}{HTML}{F7F4EE}
\definecolor{SoftBlue}{HTML}{EEF6F8}

\hypersetup{
  colorlinks=true,
  linkcolor=HeadingBlue,
  urlcolor=AccentTeal,
  pdftitle={IB Geography SL Study Pack - 2026-04-24},
  pdfauthor={IB Geography SL Preparation}
}

\color{TextInk}
\setlength{\parskip}{0.52em}
\setlength{\parindent}{0pt}
\setlength{\tabcolsep}{6pt}
\renewcommand{\arraystretch}{1.22}
\setlist[itemize]{leftmargin=1.35em,itemsep=3pt,topsep=4pt}
\setlist[enumerate]{leftmargin=1.5em,itemsep=3pt,topsep=4pt}

\newcolumntype{Y}{>{\raggedright\arraybackslash}X}
\newcolumntype{L}[1]{>{\raggedright\arraybackslash}p{#1}}

\titleformat{\section}
  {\normalfont\huge\sffamily\bfseries\color{HeadingBlue}}
  {}{0pt}{}
  [\vspace{0.15em}{\color{AccentGold}\titlerule[1pt]}]

\titleformat{\subsection}
  {\normalfont\Large\sffamily\bfseries\color{AccentTeal}}
  {}{0pt}{}

\titleformat{\subsubsection}
  {\normalfont\large\sffamily\bfseries\color{HeadingBlue}}
  {}{0pt}{}

\titlespacing*{\section}{0pt}{1.4em}{0.65em}
\titlespacing*{\subsection}{0pt}{1.0em}{0.35em}
\titlespacing*{\subsubsection}{0pt}{0.8em}{0.25em}

\pagestyle{fancy}
\setlength{\headheight}{14pt}
\fancyhf{}
\fancyhead[L]{\sffamily\color{AccentTeal}Fri 2026-04-24}
\fancyhead[R]{\sffamily\color{HeadingBlue}IB Geography SL}
\fancyfoot[C]{\sffamily\color{HeadingBlue}\thepage\ of \pageref{LastPage}}
\renewcommand{\headrulewidth}{0.4pt}
\renewcommand{\footrulewidth}{0pt}

\newcommand{\term}[1]{\textcolor{HeadingBlue}{\textbf{#1}}}
\newcommand{\exam}[1]{\textcolor{AccentTeal}{\textbf{#1}}}

\begin{document}

\begin{center}
  {\sffamily\bfseries\color{HeadingBlue}\LARGE IB Geography SL Study Pack}\\[0.25em]
  {\sffamily\bfseries\color{AccentTeal}\Large Friday 2026-04-24}\\[0.35em]
  {\sffamily\color{TextInk}Urban Environments + Population Distribution and Changing Population}
\end{center}

\vspace{0.5em}

\noindent\colorbox{SoftBlue}{%
  \begin{minipage}{0.965\textwidth}
  \sffamily
  \textbf{\textcolor{HeadingBlue}{How this pack follows the plan.}}
  The April 24 schedule asks for \term{Audio 1}: an overview and subtopic map for
  \term{Urban environments}; \term{Audio 2}: an overview of
  \term{Population distribution -- changing population}; and a \term{25 minute}
  practice block in which you write ten flashcards or a one-page summary from
  memory. This LaTeX pack is built for direct PDF conversion with textbook-style
  headings, color accents, and printable tables.
  \end{minipage}
}

\section{Today's Target}

By the end of this session, you should be able to:

\begin{itemize}
  \item name the four main subtopic clusters inside \term{Urban environments}
  \item explain how urban places change through migration, urbanization,
  suburbanization, reurbanization, gentrification, and deindustrialization
  \item distinguish \term{urban environmental stress} from \term{urban social stress}
  \item summarize the main ideas in \term{Population distribution -- changing population}
  \item interpret basic population evidence, especially maps, rates, and age-sex
  structures
  \item produce ten memory-first flashcards or one concise one-page summary
\end{itemize}

\subsection{Session Structure}

\rowcolors{2}{TableStripe}{white}
\begin{tabularx}{\textwidth}{L{0.18\textwidth}Y L{0.19\textwidth}}
\toprule
\textbf{Block} & \textbf{Focus} & \textbf{Output} \\
\midrule
Audio 1 & Urban environments overview and subtopic map & Four-cluster mental map \\
Audio 2 & Population distribution and changing population overview & Population process map \\
25 min practice & Closed-book recall, then flashcards or one-page summary & Ten cards or one page \\
\bottomrule
\end{tabularx}
\rowcolors{2}{}{}

\subsection{Fast Course Links}

\begin{tabularx}{\textwidth}{L{0.22\textwidth}Y}
\toprule
\textbf{April 24 Theme} & \textbf{Why it matters for the exam} \\
\midrule
Urban environments & This is one of the two Paper 1 options. You need both
structured-response knowledge and extended-answer evaluation. \\
Population & This is one of the three compulsory Paper 2 core units. You cannot
avoid it, so basic explanations and data skills must stay active. \\
\bottomrule
\end{tabularx}

\section{Audio 1 Script}
\subsection{Urban Environments Overview And Subtopic Map}

\textit{Target length: about 12 minutes at a natural revision pace. Read slowly
enough to pause after key terms and to sketch the four-cluster map.}

Friday, 2026-04-24. Audio 1.
\textbf{Minute 0--2: key terms and the four-cluster map.}
Today starts the first real content day after the baseline mastery board.
The focus is \term{Urban environments}, one of your two Paper 1 options.
The goal is not to memorize every city example at once. The goal is to build
a clear map of the option so that every later detail has a place to go.

Use four large headings. First, \term{variety of urban environments and urban
spatial patterns}. Second, \term{changing urban systems}. Third,
\term{urban environmental and social stresses}. Fourth, \term{building
sustainable urban systems}. Those four headings are the skeleton of the topic.
If you can recall them quickly, you can usually locate any urban question inside
one part of the course instead of treating it as a surprise.

Start with the idea of an urban environment. A city is not just a large
settlement. It is a system of land uses, flows, functions, inequalities, and
decisions. People move through it. Goods move through it. Money, information,
waste, water, energy, and political power move through it. Geography asks you
to notice how those flows produce spatial patterns. Some activities cluster
near the centre because accessibility and land value are high. Some households
move outward because they want more space or cheaper housing. Some groups are
pushed into lower-quality environments because of income, discrimination, weak
planning, or limited housing supply. A strong answer sees the city as a system,
not just as a place with buildings.

The first cluster is \term{variety and spatial pattern}. Urban environments are
not all the same. A megacity, a smaller regional city, a suburb, an informal
settlement, a commuter belt, a central business district, and a post-industrial
inner city all have different functions and problems. A rapidly growing city in
a lower-income country may be struggling to provide housing, sanitation,
transport, and formal employment for new arrivals. A wealthy city may have
better infrastructure but still face housing inequality, congestion, high
resource demand, and social segregation. A shrinking industrial city may have
vacant land, unemployment, and a smaller tax base. The word \term{urban} hides
many different realities, and IB questions often reward students who can make
that variety visible.

Spatial pattern means arrangement. It asks where different urban functions are
located and why. You should be comfortable talking about the \term{central
business district}, high land values, transport routes, residential zones,
industrial land, retail clusters, green space, edge cities, and informal
settlements. Do not treat these as static textbook diagrams only. They are
evidence of competition for space. The most accessible land often becomes the
most expensive. Lower-income groups may be pushed to peripheral areas or into
informal housing. New transport infrastructure can change accessibility and
therefore reshape land use. The basic question is always: what is located where,
and what process explains that pattern?

\textbf{Minute 2--6: processes, patterns, and causes.}

The second cluster is \term{changing urban systems}. This is where the option
becomes dynamic. The key processes include \term{urbanization}, \term{migration},
\term{suburbanization}, \term{reurbanization}, \term{gentrification}, and
\term{deindustrialization}. Urbanization means an increasing proportion of a
population living in urban areas. It can be driven by rural-to-urban migration,
natural increase within cities, or the reclassification and growth of
settlements. Migration matters because it changes city size, labor supply,
culture, housing demand, and service pressure. Push factors may include rural
poverty, conflict, environmental stress, land shortage, or lack of services.
Pull factors may include jobs, education, healthcare, security, and social
networks.

Suburbanization is the movement of people, services, or businesses from inner
urban areas to the outskirts. It is often linked with car ownership, road
building, demand for larger homes, lower land costs, and perceptions of better
quality of life. But it can create sprawl, longer commutes, loss of farmland,
and higher energy use. Reurbanization is movement or investment back into
central or inner-city areas. It may bring new housing, cultural activity,
service jobs, and higher density. But it can also raise rents and change the
social mix of neighborhoods.

Gentrification is especially useful because it connects process, pattern,
stakeholders, and evaluation. It usually involves reinvestment in a lower-income
or formerly disinvested urban area, leading to rising property values, new
businesses, physical improvement, and often displacement or exclusion of earlier
residents. The student case-study sheet mentions \term{gentrification in
Detroit}. Use that as an anchor, but make sure your class notes supply the
specific neighborhood, time period, and data your teacher expects. The exam
logic is clear even before the statistics: reinvestment can improve buildings,
services, image, and tax revenue, but it can also make the city less affordable
for people with weaker purchasing power. A balanced answer should say both.

\textbf{Minute 6--9: named example and case-study thinking.}
Treat Detroit as today's worked example for urban change, but keep the details
honest. Detroit is useful because it lets you connect deindustrialization,
decline, reinvestment, and gentrification in one place-based answer. A safe
exam chain is: manufacturing decline and population loss weakened parts of the
urban economy; later reinvestment in selected central or inner-city districts
improved buildings, services, image, and business activity; however, those
benefits were uneven because rising land values and changing neighborhood
identity can exclude lower-income residents. If your class notes used a
particular Detroit district, such as Midtown, Downtown, or another named
neighborhood, put that precise name into the answer. If your notes did not,
avoid inventing a district under exam pressure. It is better to write a careful
place-based explanation with one verified detail than to add a false statistic.

Deindustrialization is the decline of manufacturing employment and industrial
activity in a place. It matters because many cities grew around factories,
ports, railways, warehouses, and working-class housing. When industry closes or
moves, the city can be left with unemployment, derelict land, lower municipal
revenue, and social stress. But deindustrialized land can also become a target
for regeneration: new housing, offices, cultural quarters, parks, or mixed-use
development. Again, the exam habit is evaluation. Regeneration may improve the
built environment, but who gains? Existing residents, new residents, developers,
tourists, or the local government?

\textbf{Minute 9--10: stresses and management responses.}
The third cluster is \term{urban stress}. Divide it into environmental stress
and social stress. Environmental stress includes air pollution, noise, traffic
congestion, waste disposal, pressure on water supply, urban heat islands, flood
risk, loss of vegetation, and high energy demand. These are not random side
effects. They often come from density, car dependence, industrial activity,
impermeable surfaces, weak regulation, and rapid growth that outruns
infrastructure. If you are asked to explain an urban environmental problem, try
to build a causal chain: growth or land-use change leads to pressure on a system,
which creates a measurable environmental impact, which then affects people
unevenly.

Social stress includes housing shortage, informal settlements, homelessness,
inequality, segregation, service gaps, crime, unemployment, displacement, and
political conflict over land use. Social stress is often linked to who has
access to urban opportunities. A city may produce wealth and still exclude large
groups from secure housing, reliable transport, good schools, healthcare, green
space, or political voice. The strongest geography answers avoid vague claims
such as "cities are crowded" and instead explain why crowding, inequality, or
service pressure occurs in a named place or type of city.

The fourth cluster is \term{sustainable urban systems}. This is the management
part of the topic. A sustainable city tries to reduce environmental damage,
improve social inclusion, and remain economically viable. Strategies can include
public transport, transit-oriented development, cycling and walking
infrastructure, mixed land use, compact development, energy-efficient buildings,
urban greening, flood-sensitive design, recycling, waste reduction, affordable
housing, slum upgrading, brownfield redevelopment, and participatory planning.
But sustainable does not mean perfect. A metro system may cut emissions but be
too expensive for some users. Urban greening may cool the city but also raise
nearby property values. Compact development may reduce sprawl but increase
pressure on local services if planning is weak.

To revise sustainability well, use three questions: sustainable for whom,
sustainable at what scale, and sustainable for how long? A policy can look
successful at city scale while harming one neighborhood. It can reduce emissions
but worsen affordability. It can improve short-term image but fail to solve
long-term inequality. IB extended answers often reward exactly this kind of
judgment.

Keep a short strategy bank in your head. Public transport can reduce congestion,
air pollution, and carbon emissions, but only if routes are affordable,
reliable, and connected to where people actually live and work. Urban greening
can reduce heat-island effects, improve mental health, absorb rainfall, and
make districts more attractive, but it can also raise nearby property values if
housing protection is weak. Brownfield redevelopment can reuse derelict land and
slow outward sprawl, but it may be expensive to clean contaminated sites and may
produce luxury housing instead of affordable housing. Slum upgrading can improve
water, sanitation, electricity, roads, and tenure security without forcing
residents away, but it needs long-term funding and community trust. In a timed
answer, one strong paragraph can compare two of these strategies and judge which
is more socially inclusive.

Now connect the four clusters. Urban variety gives you the types of places and
patterns. Changing urban systems gives you the processes that reshape those
places. Urban stress gives you the consequences of rapid or uneven change.
Sustainable urban systems gives you the responses and the evaluation. That is
the whole option in one sequence: pattern, process, pressure, response.

Use that sequence with any urban question. If the question is about land use,
start with pattern and accessibility. If it is about migration, start with
process and consequences. If it is about pollution or housing, start with stress
and unequal impact. If it is about planning, start with response and evaluation.
This prevents your answer from becoming a list.

\textbf{Minute 10--12: model command-term paragraph.}
Here is a model \exam{analyse} paragraph in spoken form. Gentrification can
create both opportunities and challenges because it changes the built
environment and the social geography of a neighborhood at the same time. In
Detroit, reinvestment in selected inner-city areas can improve derelict
buildings, attract services, create jobs, and strengthen the city's image after
industrial decline. Those are real urban benefits because they can increase
economic activity and make some districts more liveable. However, the same
process can raise rents and property values, which may exclude residents who
have less income or less secure housing. This means the benefits of regeneration
are spatially and socially uneven. The strongest judgment is not that
gentrification is simply good or bad. It is that gentrification can improve
urban environments physically while deepening social inequality unless planning,
affordable housing, and community participation protect existing residents.

If the command term is \exam{evaluate} instead of \exam{analyse}, push the
judgment further. You could say: the success of gentrification or regeneration
depends on the criteria used. If success means attracting investment and
improving the built environment, selected Detroit districts may show positive
change. If success means reducing inequality and protecting existing residents,
the judgment becomes less secure because rising values can shift the benefits
toward developers, new residents, and higher-income consumers. That is why an
evaluative answer should name the winners, name the losers, and state whether
the strategy solves the original urban problem or simply moves it elsewhere.

For the 25-minute block later, your first five flashcards should probably come
from Urban environments. Suggested card fronts are: urbanization,
suburbanization, gentrification, deindustrialization, and sustainable urban
management. For each card back, include a definition, one cause, one consequence,
and one case-study hook. That format turns a term into an exam-ready idea.

Before the pause, add one visual-stimulus habit for urban questions. If an exam
map, photograph, or graph shows an urban pattern, do not start with explanation.
Start with visible evidence. Name the pattern, such as high-density housing near
the centre, traffic congestion along a main route, new investment in an inner
district, or informal housing at the urban edge. Then give the process that
could explain it: accessibility, land value, migration, planning, income
inequality, or infrastructure shortage. Finally, state one implication for
people or management. That three-step move, evidence, process, implication, is
how you turn an urban stimulus into a geographical answer instead of a vague
description.

The same habit works for a table of urban data. If one district has higher rent,
lower green space, or longer commuting time than another, describe the contrast
first, then explain it through accessibility, income, land value, planning, or
historic investment. Finish by linking the pattern to social or environmental
stress. This is the quickest route from observation to analysis when the
stimulus is unfamiliar.

Pause now and redraw the four-cluster map from memory. Do not look back yet.
Write: variety and pattern, changing systems, stresses, sustainability. Under
each heading, add at least three terms. If one heading stays blank, that is
today's weakness to fix first.

\subsection{Urban Environments Key Terms}

\rowcolors{2}{TableStripe}{white}
\begin{tabularx}{\textwidth}{L{0.24\textwidth}Y Y}
\toprule
\textbf{Term} & \textbf{Core meaning} & \textbf{How to use it in an answer} \\
\midrule
Urbanization & Increasing share of a population living in urban areas. &
Link to migration, natural increase, jobs, services, housing demand, and
infrastructure pressure. \\
Suburbanization & Movement of people and activities outward from inner urban
areas. & Evaluate quality-of-life gains against sprawl, commuting, land loss,
and car dependence. \\
Reurbanization & Return of people or investment to inner urban areas. & Link to
regeneration, higher density, new services, and possible displacement. \\
Gentrification & Reinvestment that raises land values and changes the social mix
of a lower-income area. & Balance physical improvement against affordability and
displacement. \\
Deindustrialization & Decline of manufacturing employment and industrial land
use. & Explain unemployment, dereliction, brownfield land, and regeneration
possibilities. \\
Urban stress & Pressure on environmental or social systems in cities. & Separate
environmental impacts from social inequality and show who is most affected. \\
Sustainable urban management & Planning that aims to reduce environmental harm
and improve social and economic outcomes. & Judge the strategy's benefits,
limits, scale, and equity. \\
\bottomrule
\end{tabularx}
\rowcolors{2}{}{}

\newpage

\section{Audio 2 Script}
\subsection{Population Distribution And Changing Population Overview}

\textit{Target length: about 12 minutes at a natural revision pace. Pause after
each data-skill instruction and sketch one population pyramid from memory.}

Friday, 2026-04-24. Audio 2.
\textbf{Minute 0--2: key terms and the core map.}
The second focus today is \term{Population distribution -- changing population},
one of the compulsory Paper 2 core units. This matters because Paper 2 does not
let you avoid the core. Population can appear through structured questions,
maps, graphs, population pyramids, data tables, or an extended response link.
Your target today is a strong overview: where people are, why population changes,
how structure matters, and how policy responds.

Start with \term{population distribution}. Distribution means the way people are
spread across space. Some areas are densely populated; others are sparsely
populated. Geography asks you to explain that pattern using both physical and
human factors. Physical factors include relief, climate, soils, water supply,
natural hazards, vegetation, and access to coasts or rivers. Human factors
include employment, transport, history, trade, government policy, conflict,
technology, and service provision. A river valley, coast, or fertile plain can
attract population because it supports farming, trade, transport, and settlement.
A desert, mountain range, dense forest, or cold high-latitude environment may
remain sparsely populated because water, access, or production is limited.

But avoid one simple trap: physical geography never explains everything by
itself. Technology and wealth can change what is possible. A dry region may
support a large city if it has desalination, energy, imports, and strong
infrastructure. A fertile region may remain underpopulated because of conflict,
disease, weak transport, or political instability. So in an exam answer, use
physical factors and human factors together. That gives you a more geographical
explanation.

\textbf{Minute 2--6: processes, patterns, and causes.}
Next, population changes through \term{births}, \term{deaths}, and
\term{migration}. The basic formula is simple: population change equals natural
increase plus net migration. Natural increase is births minus deaths. Net
migration is immigration minus emigration. A country can grow because many more
people are born than die, because more people arrive than leave, or because both
happen at the same time. A country can shrink because deaths exceed births,
because more people leave than arrive, or because low fertility combines with
out-migration. Keep this formula active. It helps you interpret almost any
population graph.

Now focus on \term{fertility}. Fertility refers to births, often measured by the
total fertility rate. Fertility is shaped by income, education, gender equality,
urbanization, child survival, contraception, religion, culture, housing cost,
state policy, and the economic role of children. In many places, fertility falls
as child survival improves, women gain more education, urban living becomes more
expensive, and families choose to have fewer children. But there is no automatic
single path. Government policy, culture, inequality, and access to healthcare can
make fertility change faster or slower.

\term{Mortality} refers to deaths. Mortality is shaped by healthcare, nutrition,
sanitation, disease, conflict, disasters, age structure, and living conditions.
Falling death rates often reflect improved public health, clean water, vaccines,
medical care, food supply, and education. But mortality can rise or remain high
where conflict, disease, poverty, environmental hazards, or weak healthcare
increase risk. In a short answer, do not just say "better healthcare." Add the
mechanism: better healthcare reduces deaths from disease or childbirth, which
raises life expectancy and alters the age structure.

Migration is the third driver of population change. It can be internal or
international, voluntary or forced, temporary or permanent. Push factors include
unemployment, low wages, conflict, persecution, poor services, land shortage,
environmental stress, or lack of opportunity. Pull factors include jobs,
education, healthcare, safety, family networks, political freedom, and higher
wages. Migration changes both sending and receiving areas. Sending areas may
lose young workers but receive remittances. Receiving areas may gain labor and
cultural diversity but face housing, service, or political pressure. As with
urban geography, the strongest answers evaluate who benefits and who loses.

Population structure is about the composition of a population by age and sex.
The main visual tool is the \term{population pyramid}. A wide base usually
suggests high birth rates and a youthful population. A narrow base with a large
middle and upper section suggests low birth rates, ageing, and longer life
expectancy. A pyramid with a working-age bulge may suggest potential for a
demographic dividend if education, jobs, and governance are strong enough to use
that labor force productively. But if employment is weak, the same youthful
structure can produce pressure on schools, housing, services, and political
stability.

When reading a population pyramid, use four steps. First, describe the shape.
Second, infer the process that created it. Third, state the consequence. Fourth,
link to policy. For example: a wide base suggests high fertility. High fertility
creates a youthful age structure. That increases demand for schools, childcare,
healthcare, and future jobs. Policy may therefore focus on education, maternal
health, family planning, or employment creation. This sequence turns visual
description into analysis.

\textbf{Minute 6--9: named examples and case-study thinking.}

The opposite structure is ageing. An ageing population has a growing share of
older people. It is often caused by low fertility and long life expectancy. The
student case-study sheet mentions \term{Japan's ageing population}; this is a
useful anchor because Japan is commonly used to show the pressures of a high
elderly share, long life expectancy, low fertility, labor shortages, pension
pressure, healthcare demand, and rural service decline. In an answer, the core
chain is: low fertility plus long life expectancy creates a high old-age
dependency ratio, which can increase the tax burden on workers and pressure
public services. But do not present older people only as a problem. Ageing can
also support new markets, community roles, and technological innovation in
healthcare and robotics. A more balanced answer is usually stronger.

Youthful populations create a different set of opportunities and pressures. A
high proportion of young people can mean high demand for schools, vaccination,
housing, water, sanitation, and future jobs. If the economy expands and education
is strong, a youthful population can become a demographic dividend: a large
working-age population that supports growth. If employment and governance are
weak, the same structure can intensify unemployment, informal work, migration,
or unrest. The age structure itself is not destiny; outcomes depend on policy
and opportunity.

Two terms deserve special attention here: \term{dependency ratio} and
\term{demographic dividend}. Dependency ratio matters because it turns age
structure into economic pressure. A high youth dependency ratio points toward
pressure on schools, childcare, vaccination, and future job creation. A high
old-age dependency ratio points toward pressure on pensions, healthcare, social
care, and the tax base. Demographic dividend is the positive possibility: a
large working-age share can support growth, savings, and productivity. But it is
not automatic. It needs education, health, job creation, gender inclusion, and
stable governance. Without those conditions, a youthful population may become
underemployment and migration pressure rather than a dividend.

Population policy is the management side of the unit. States may try to reduce
fertility, raise fertility, manage migration, support ageing populations, or
redistribute population. The student case-study sheet mentions \term{China's
one-child policy}. Use it as an example of an antinatalist policy designed to
slow population growth. The exam value is not just knowing that the policy
existed. You need cause, method, effect, and evaluation. It reduced births and
slowed growth, but it also contributed to ageing, a smaller future workforce,
and gender imbalance. A careful answer recognizes that population policy can
solve one pressure while creating another.

The student case-study sheet also mentions \term{Singapore's pronatalist and
antinatalist policies}. This is valuable because it shows policy can change
direction as demographic needs change. Singapore used antinatalist measures when
rapid growth was viewed as a problem, then later shifted toward pronatalist
measures when low fertility, ageing, and labor-supply concerns became more
important. Use this case to show evaluation: incentives such as baby bonuses,
childcare support, parental leave, or housing benefits may help, but they do not
automatically overcome high living costs, career pressure, gender expectations,
and delayed marriage.

\textbf{Case-study comparison.}
Use the three population anchors as a contrast set, not as three isolated facts.
Japan helps you explain what happens when low fertility and high life expectancy
produce an older age structure. China helps you evaluate a restrictive
antinatalist policy: it reduced one pressure, rapid growth, but strengthened
other pressures, especially ageing and gender imbalance. Singapore helps you
show that population policy changes over time: when rapid growth is the problem,
policy may discourage births; when low fertility and ageing become the problem,
policy may encourage births. In a Paper 2 answer, that comparison is valuable
because it lets you move from description to judgment. Population policy is
rarely a clean success or failure. It works within culture, housing cost, labor
markets, gender expectations, healthcare, and migration rules.

\textbf{Minute 9--10: visual and data-skill practice.}
Now connect population to Paper 2 skills. The IB assessment model values the
ability to interpret patterns and processes in unfamiliar information. That
means you must practise moving from data to explanation. If a map shows dense
population along a coast, describe the pattern and then explain possible access,
trade, climate, flat land, or historical settlement factors. If a line graph
shows falling fertility, explain likely social and economic changes. If a table
shows ageing, link it to healthcare, pensions, labor supply, dependency ratio,
and possible migration policy. Data earns marks only when you turn it into
geographical meaning.

Try a spoken data drill. Imagine Population A has 36 percent aged 0--14, 59
percent aged 15--64, and 5 percent aged 65 or over. Imagine Population B has 12
percent aged 0--14, 59 percent aged 15--64, and 29 percent aged 65 or over.
The first move is description: Population A is more youthful, while Population B
is much older. The second move is explanation: Population A probably has higher
fertility, while Population B probably has lower fertility and longer life
expectancy. The third move is consequence: Population A needs schools and future
jobs, while Population B needs pensions, healthcare, and labor-force planning.
The fourth move is policy: Population A may invest in education and job creation,
while Population B may encourage higher fertility, later retirement, automation,
or immigration. That is the whole Paper 2 data skill in miniature.

A useful memory frame for this whole unit is \term{pattern, process, structure,
policy}. Pattern asks where people are. Process asks how births, deaths, and
migration change numbers. Structure asks what the age-sex composition looks like.
Policy asks how governments respond. If you can recall those four words, you can
organize almost any population answer.

Also link population back to yesterday's course map. Population connects to the
other core units. Rapid population growth can increase water, food, energy,
housing, and transport demand. Ageing can change consumption patterns, tax
systems, and healthcare needs. Migration can be shaped by climate stress or
resource insecurity. Population is therefore not an isolated unit. It is the
human pressure and human capacity side of global change.

\textbf{Minute 10--12: model command-term paragraph.}
Here is a model \exam{examine} paragraph in spoken form. China's one-child
policy shows that a population policy can be effective in one aim while creating
serious long-term costs. The policy was designed to slow rapid population
growth by limiting births, so it addressed pressure on food, housing, education,
and employment. In that narrow sense, it can be seen as effective because it
helped reduce fertility and slow natural increase. However, its consequences
also show why population policy needs evaluation. Fewer births contributed to
an ageing population and a higher old-age dependency burden. The policy also
contributed to gender imbalance and raised ethical concerns about state control
over family size. A balanced judgment is therefore that the policy reduced one
demographic pressure but created new social and economic pressures that later
policy had to manage.

For today's flashcards, choose at least five population terms. Good card fronts
are: population distribution, natural increase, fertility rate, migration,
population pyramid, ageing population, dependency ratio, demographic dividend,
and population policy. On the back of each card, write a definition, one cause,
one consequence, and one example. If you use a case-study card, make the back
specific: place, policy or issue, intended effect, actual effect, and one
evaluative limitation.

One final exam reminder: keep scale visible. A population trend can be national,
regional, urban, rural, household-level, or age-group specific. A country may
look demographically stable overall while one region is losing young adults and
another city is gaining migrants. When you write, say which scale you mean.

That scale habit prevents overgeneralization, which is a common weakness in
Paper 2 answers. Precise scale also makes examples easier to evaluate because
impacts are not shared equally.

Pause now and answer from memory. What are the three drivers of population
change? What are the four steps for reading a population pyramid? What is one
benefit and one problem linked to an ageing population? What is one strength and
one limitation of a population policy you have studied? If you can answer those
clearly, you have the overview needed for today.

\subsection{Population Key Terms}

\rowcolors{2}{TableStripe}{white}
\begin{tabularx}{\textwidth}{L{0.24\textwidth}Y Y}
\toprule
\textbf{Term} & \textbf{Core meaning} & \textbf{How to use it in an answer} \\
\midrule
Population distribution & The spread of people across space. & Explain using
physical and human factors together. \\
Population density & Number of people per unit area. & Use carefully because
average density can hide internal variation. \\
Natural increase & Births minus deaths. & Link to fertility, mortality, growth,
and age structure. \\
Net migration & Immigration minus emigration. & Explain impacts on both sending
and receiving areas. \\
Population pyramid & Graph showing age and sex structure. & Describe shape,
infer process, state consequence, and link to policy. \\
Dependency ratio & Relationship between dependent age groups and working-age
population. & Use for ageing, youthfulness, tax pressure, and service demand. \\
Demographic dividend & Possible economic gain from a large working-age
population. & Evaluate conditions needed: education, jobs, health, governance. \\
Population policy & Government action to influence fertility, mortality,
migration, or distribution. & Judge effectiveness and unintended consequences. \\
\bottomrule
\end{tabularx}
\rowcolors{2}{}{}

\newpage

\section{25-Minute Practice Block}

\subsection{Minute 0-5: Closed-Book Brain Dump}

Without checking the scripts, make two headings:
\term{Urban environments} and \term{Population}. Under each heading, write as
many terms, processes, examples, and uncertainties as you can. Use fragments,
not full sentences. Speed matters more than neatness.

\begin{center}
\renewcommand{\arraystretch}{1.45}
\begin{tabularx}{0.96\textwidth}{Y Y}
\toprule
\textbf{Urban environments} & \textbf{Population} \\
\midrule
\rule{0pt}{4.4cm} & \rule{0pt}{4.4cm} \\
\bottomrule
\end{tabularx}
\end{center}

\subsection{Minute 5-18: Choose One Output}

Choose \term{Option A} or \term{Option B}. Do not do both unless you have extra
time.

\noindent\colorbox{SoftTint}{%
  \begin{minipage}{0.965\textwidth}
  \sffamily
  \textbf{\textcolor{HeadingBlue}{Option A: Ten flashcards.}}
  Write ten cards from memory. Use five Urban environments cards and five
  Population cards. Each card back should contain:
  \textbf{definition + cause/process + consequence + example hook}.
  \end{minipage}
}

\vspace{0.5em}

\noindent\colorbox{SoftBlue}{%
  \begin{minipage}{0.965\textwidth}
  \sffamily
  \textbf{\textcolor{HeadingBlue}{Option B: One-page summary.}}
  Fit the whole day onto one page. Use four boxes:
  \textbf{urban pattern}, \textbf{urban change}, \textbf{population process},
  and \textbf{population policy}. Each box needs terms, one example, and one
  evaluative point.
  \end{minipage}
}

\subsection{Minute 18-25: Check And Fix}

Use the checklist below. Add corrections in a different color if writing by
hand.

\begin{itemize}
  \item Did every flashcard or summary box include an example hook?
  \item Did at least two entries include evaluation, not just description?
  \item Did you separate environmental urban stress from social urban stress?
  \item Did you include births, deaths, and migration as population drivers?
  \item Did you include at least one case-study anchor from your own class notes?
\end{itemize}

\section{Case-Study Anchors To Fill Precisely}

The current student case-study sheet lists the following anchors. Keep them, but
add the exact facts, dates, locations, and statistics your class has used.

\rowcolors{2}{TableStripe}{white}
\begin{tabularx}{\textwidth}{L{0.23\textwidth}Y Y}
\toprule
\textbf{Topic} & \textbf{Anchor already listed} & \textbf{What to add before the exam} \\
\midrule
Urban environments & Gentrification in Detroit & Neighborhood or project name,
evidence of reinvestment, who benefits, who is displaced or excluded, and one
judgment. \\
Population policy & China's one-child policy & Dates, methods, fertility or
growth impact, ageing impact, gender imbalance, and one evaluation. \\
Ageing population & Japan's ageing population & Evidence of ageing, dependency
pressure, labor or healthcare impacts, government responses, and one limitation. \\
Population policy & Singapore's pronatalist and antinatalist policies & Why the
policy direction changed, examples of incentives or restrictions, outcomes, and
one limitation. \\
\bottomrule
\end{tabularx}
\rowcolors{2}{}{}

\section{Exam-Style Mini Practice}

\subsection{Short Structured Questions}

Answer any two in 5 minutes each.

\begin{enumerate}
  \item \exam{Describe} two features of an urban land-use pattern.
  \item \exam{Explain} one reason why suburbanization may create environmental
  stress.
  \item \exam{Explain} how migration can change the population structure of a
  receiving area.
  \item \exam{Suggest} one reason why an ageing population can create economic
  pressure.
\end{enumerate}

\subsection{One Paragraph Challenge}

Choose one prompt and write 8--10 lines.

\begin{enumerate}
  \item \exam{Analyse} how gentrification can create both benefits and costs in
  an urban area you have studied.
  \item \exam{Examine} one population policy and its unintended consequences.
\end{enumerate}

\subsection{Paragraph Frame}

\noindent\fcolorbox{RuleGray}{SoftTint}{%
  \begin{minipage}{0.965\textwidth}
  \textbf{\textcolor{HeadingBlue}{Claim:}} state the process or policy clearly.\\
  \textbf{\textcolor{HeadingBlue}{Evidence:}} add a place, data point, or class
  case-study detail.\\
  \textbf{\textcolor{HeadingBlue}{Explanation:}} show the causal chain.\\
  \textbf{\textcolor{HeadingBlue}{Evaluation:}} identify who benefits, who loses,
  or what the limitation is.
  \end{minipage}
}

\section{End-Of-Day Output}

You are finished with April 24 if you have:

\begin{itemize}
  \item listened to or read both audio scripts
  \item redrawn the Urban environments four-cluster map from memory
  \item summarized Population using pattern, process, structure, and policy
  \item produced ten flashcards or one one-page summary
  \item updated yesterday's 1--5 mastery rating for Urban environments and
  Population
  \item wrote at least one error-log note: the mistake, the correction, and the
  next action
  \item marked the weakest three items to carry into the April 25 Urban
  environments depth session
\end{itemize}

\end{document}
